\section{Related Work}
\label{sec:related}

Counterfactual protocol evaluation intersects four literatures: learning from
incomplete temporal records, measurement design under an assumed process
model, evaluation across acquisition regimes, and statistical identification
of observation schemes. The main distinctions are whether the measurements
required by an alternative occur in the observed data and whether its value is
computable under a specified model or must be identified from the realised
measurement--target law.

\subsection{Learning from Incomplete Temporal Trajectories}
\label{sec:rw-incomplete}

Functional data analysis estimates population mean and covariance functions
from sparse or irregular measurements under sampling and smoothness
assumptions
\citep{ramsay2005functional,yao2005sparse,hall2006properties,li2010uniform}.
Longitudinal methods also allow observation times to depend on latent or
previously observed outcomes
\citep{pullenayegum2016irregular,weaver2023informative}. These estimators pool
within-subject pairs across units, so recovery of a covariance region requires
the sampling design to supply informative pairs there, together with the
assumptions imposed by the estimator. Under a single fixed protocol, times or
supports unique to an alternative may occur for no unit. The corresponding
covariance blocks are then not determined by the measurements alone; in the
linear-Gaussian analysis, the aggregate target adds the cross-covariance and
variance constraints characterised in Section~\ref{sec:identifiability}.

Machine-learning methods for irregular time series instead learn prediction or
interpolation rules directly from incomplete records. GRU-D models missingness
indicators and elapsed times, whereas latent ODEs and neural controlled
differential equations represent irregular observations in continuous time
\citep{che2018grud,rubanova2019latent,kidger2020neural,shukla2020survey}.
Early-prediction and active feature acquisition methods determine when an
observed prefix is sufficient or which feature to acquire next
\citep{xing2009early,shim2018joint}. Multiple-instance and label-proportion
methods also use aggregate labels, but retain a fixed observation scheme
\citep{dietterich1997multiple,quadrianto2009labels,zhou2018brief}. These methods
operate on observation patterns represented in the data. By contrast, an
alternative protocol here may require measurements absent for every unit in
the realised sample; its joint law with the target is therefore not available
as a marginal of the realised law.

\subsection{Measurement Design under an Assumed Model}
\label{sec:rw-design}

Classical optimal design allocates measurements to estimate parameters or
linear functionals efficiently under a statistical model
\citep{elfving1952optimum,pukelsheim2006optimal}. Bayesian and goal-oriented
design optimise expected information or posterior uncertainty about a
quantity of interest under a model or prior
\citep{lindley1956information,chaloner1995bayesian,attia2018goal,
alexanderian2021optimal,rainforth2024modern}. Active learning, sensor
selection, Gaussian-process sensor placement, and kernel quadrature likewise
optimise uncertainty, information gain, estimation error, or integration
accuracy under a fitted model, covariance, or kernel
\citep{mackay1992information,settles2012active,joshi2009sensor,
krause2008near,bach2017equivalence,briol2019probabilistic}.

Optimal design for longitudinal and functional data is especially close to the
present temporal setting: pilot data can estimate dependence and guide future
observation times, including designs for predicting a scalar response
\citep{ji2017optimal}. Such procedures ask which measurements to collect once
a model makes the criterion computable. Counterfactual protocol evaluation
asks the preceding question: whether the value in
Definition~\ref{def:values} is determined by the realised law at all. Once the
required dependence has been identified or calibrated, Section~\ref{sec:design}
becomes a target-aware sensor-selection problem. Its objective is the
population $R^2$ of the Bayes-optimal predictor of a fixed temporal aggregate,
not trajectory-reconstruction error or target-free information gain. For
noiseless point measurements and a linear aggregate it coincides with a
kernel-quadrature criterion; measurement noise and nonlinear aggregates move
beyond that case.

\subsection{Evaluation across Acquisition Regimes}
\label{sec:rw-evaluation}

Active feature acquisition performance evaluation (AFAPE) evaluates a new
acquisition policy from retrospective data generated by another acquisition
process \citep{vonkleist2025evaluation}. Depending on its assumptions, AFAPE
uses missing-data or offline-policy-evaluation methods, including direct,
inverse-probability-weighted, and doubly robust estimators. More generally,
off-policy evaluation requires overlap or related support conditions between
the behaviour and target regimes
\citep{precup2000eligibility,kallus2020double,sachdeva2020deficient}.

A fixed observation protocol can be represented as a deterministic acquisition
policy, but an alternative considered here may assign positive demand to
measurements having zero acquisition probability under the realised protocol.
Reweighting observed acquisition trajectories cannot identify their
distribution or predictive value. Identification must instead come from
restrictions linking observed and unobserved components of the same latent
trajectory to the aggregate target. Theorem~\ref{thm:minimal} and
Theorem~\ref{thm:impossibility} show that these links need not determine the
alternative's value, whereas Proposition~\ref{prop:augmentation} gives
conditions under which additional measurements remove the relevant ambiguity.
Graphical missing-data theory provides a related account of recoverability from
systematically incomplete observations
\citep{mohan2021graphical,nabi2020full}. Here, however, the acquisition rule is
fixed and known; the unresolved object is the latent temporal dependence
between measurements that were and were not collected.

\subsection{Identification, Comparison, and Predictive Limits of Observation Schemes}
\label{sec:rw-identification}

Statistical identification asks whether a functional is constant over all
latent models inducing the same observed law; partial identification retains
the resulting set of compatible values
\citep{koopmans1950identification,rothenberg1971identification,
manski2003partial,tamer2010partial}. We apply this principle to the predictive
value of an undeployed protocol. Definition~\ref{def:invisible} exposes the
observational-equivalence geometry in the linear-Gaussian model,
Theorem~\ref{thm:impossibility} gives a local non-identification certificate,
and Proposition~\ref{prop:permutation} gives an exact construction for
nonlinear aggregates. Proposition~\ref{prop:augmentation} conversely shows that
a specified protocol value can be identified without recovering the complete
latent covariance.

Bayes-error analyses characterise predictive limits for a fixed observed
feature--target law \citep{fukunaga1987bayes,ishida2023performance}. The present
question is whether that law determines the predictive limit under a different
feature-generating protocol. Blackwell's comparison of experiments and its
extensions order observation schemes across classes of decision problems
\citep{blackwell1953equivalent,lecam1990asymptotics,torgersen1991comparison};
protocol value is instead a scalar criterion for one fixed target and loss.
Transportability and data-fusion methods infer quantities outside a realised
regime through causal or invariance assumptions across populations or data
sources \citep{pearl2014external,bareinboim2016fusion}. Here the population and
target remain fixed, and the regime change concerns which measurements of the
same latent trajectory are collected.

The Gaussian identity used to evaluate a fixed protocol in
Proposition~\ref{prop:risk-identity} was derived by
\citet{zhang2026label}. The present work studies the identification of an
undeployed protocol's value, measurement augmentation that restores it, and
finite-calibration comparison and design.
